
\section{Outcomes}

One must approach all endeavours with a \textbf{set of goals}. Here are the expected outcomes of a university education, the chemistry program, and this course.

\subsection{University Outcomes}

Students who have completed a \textbf{bachelor of science degree} will have achieved the following goals:

\begin{itemize}
\item Be able to \textbf{gather information} through a well documented process of \textbf{inquiry}.
\item Be able to \textbf{analyze and interpret} that information and produce a \textbf{conclusion}.
\item Be able to effectively \textbf{communicate} the conclusion, the original data and the methods used in the analysis.
\item Have mastered \textbf{tools and methods} to accomplish the above within their field of endeavour.
\item Be able to work with others demonstrating \textbf{leadership and collaborative skills} and to function as a member of an interdisciplinary \textbf{problem solving} team.
\item Have become \textbf{independent thinkers} who are responsible for their own learning.
\item Have demonstrated the principles of \textbf{academic integrity} in all their scholarly activities.
\end{itemize}

\subsection{Program Outcomes}

Students who complete a \textbf{BSc in Chemistry} will have achieved the following goals:

\begin{itemize}
\item Have developed an understanding of the principles of various \textbf{fields of chemistry} (organic, inorganic, physical, analytical, and biochemistry).
\item Be able to \textbf{safely conduct oneself in a chemical laboratory} and any location where chemistry is applied.
\item Be able to \textbf{design and carry out scientific experiments}, accurately \textbf{document} methods and \textbf{analyze} the results of such experiments.
\item Be able to \textbf{use a variety of modern instruments} for separating and analyzing mixtures and for characterizing molecules and materials and to \textbf{apply these skills} to appropriate chemical problems.
\item Be able to \textbf{interpret} the modern primary \textbf{literature} of chemistry and to \textbf{communicate} the results, conclusions, and relevance of scientific experiments to a specific audience in a variety of formats, including \textbf{peer-review} publications.
\item Understand the \textbf{role of chemistry} in making the world around us better for everyone and understand the risks and harms that can be caused by chemistry.
\item Be able to \textbf{explore new areas of research} in both chemistry and allied fields of science and technology.
\item \textbf{Honours students} will have greater exposure to upper-level chemistry courses and research under a supervisor with the goal of producing a research \textbf{thesis} and publicly presenting the results.
\end{itemize}

\subsection{Course Outcomes}

Students who successfully complete this course of study in \textbf{physical organic chemistry} will have achieved the following goals:

\begin{itemize}
\item Be able to \textbf{describe molecular structure} by the use of valence bond theory and molecular orbital theory.
\item Be able to use \textbf{computational tools} to determine optimized structures of small molecules and interpret the energetics of simple reactions.
\item Be able to \textbf{design experiments} to determine the kinetic and thermodynamic properties of a reaction and to \textbf{interpret the results}.
\item Be familiar with different types of \textbf{reaction mechanisms} and be able to describe different types of \textbf{reactive intermediates}.
\item Be able to \textbf{propose a reaction mechanism} after examining the reactants and products of a reaction.
\item Be able to use reaction kinetics as a tool for \textbf{understanding reaction mechanisms}.
\item Develop a deeper understanding of \textbf{acid/base equilibria} in aqueous and non-aqueous systems and determine whether a reaction is \textbf{acid or base catalyzed}.
\item Be able to apply fundamental concepts of \textbf{chemical and biochemical catalysis}.
\item Be able to \textbf{investigate} a hypothetical reaction mechanism by:
\begin{itemize}
\item using the \textbf{Hammett equation} as a tool in studies of organic reactions;
\item using other \textbf{linear free energy relationships} (LFERs);
\item incorporating \textbf{isotopes};
\item applying a variety of \textbf{other methods}.
\end{itemize}
\item Be able to explain \textbf{steric and electronic effects} in several important reactions.
\item Develop an understanding of the \textbf{transition state} and how its structure changes with steric and electronic effects.
\item Be better able to choose suitable \textbf{solvents and conditions} for a reaction.
\item Understand how \textbf{parameter systems} developed using physical organic chemistry are used in \textbf{quantitative structure--activity relationships (QSAR)} for drug design.
\item Have improved \textbf{computer skills} in data analysis, chemical investigation, and record keeping.
\end{itemize}

\subsection{Self Evaluation}

Almost all students in this course will be in the final year of the chemistry program. \textbf{Now is the time} to evaluate yourself as you consider the goals for \textbf{university and program outcomes}. How do you feel? Do you have a sense that you are well on the way to achieving them? Are there any weaknesses that could be addressed while taking this course?

Also consider the goals for the \textbf{course outcomes}. You might feel that you have accomplished some of them already. \textbf{Return to this list} every few weeks and check off the goals you feel you have reached or are progressing toward. By monitoring your progress, you will be better able to see where more work is needed to ensure that all goals are met. Never hesitate to \textbf{ask for more} in one area or another.


